\setcounter{section}{-1}
\section{一线三等角：快速回顾}
先看一幅熟悉的图：\(A\)、\(B\)、\(C\) 在同一直线上，\(B\) 在 \(A\)、\(C\) 之间；\(D\)、\(E\) 在直线 \(AC\) 的同侧，且
\[
\angle DAB=\angle DBE=\angle BCE=90^\circ,\qquad BD=BE.
\]
\diagram{\diReview}

直线上的三个角满足
\[
\angle DBA+\angle DBE+\angle EBC=180^\circ,
\]
所以 \(\angle DBA+\angle EBC=90^\circ\)。而在直角三角形 \(DAB\) 中，
\(\angle BDA+\angle DBA=90^\circ\)，故 \(\angle BDA=\angle EBC\)。

再结合 \(\angle DAB=\angle BCE\) 和 \(BD=BE\)，得到
\[
\triangle DBA\cong\triangle BEC\quad(\mathrm{AAS}),
\]
于是 \(\boldsymbol{AB=CE}\)、\(\boldsymbol{AD=BC}\)。注意对应顺序：
\(D\leftrightarrow B\)、\(B\leftrightarrow E\)、\(A\leftrightarrow C\)。

\begin{strand}
\textbf{识别顺序}\quad 一条直线 \(\longrightarrow\) 三个直角
\(\longrightarrow\) 等斜边 \(\longrightarrow\) 全等后交换对应直角边。

本章只回顾这一种直角全等型，下面“作双垂”的证明会连续用到它两次。
\end{strand}
\clearpage
